\twocolumn[{
\section{Results: Stopping Reasons}
\begin{center}\small
\begin{tabular}{lrrrrrrr}
\toprule
Model & Sample F1 & Macro F1 & Safety & Funding & Efficacy & Recruit. & Other \\
\midrule
Supervised LR & 0.8782 & 0.6990 & 0.8571 & 0.2857 & 0.8125 & 0.9796 & 0.5600 \\
Initial K-Means + LR & 0.8547 & 0.6690 & 0.8235 & 0.4000 & 0.7407 & 0.9524 & 0.4286 \\
Retained K-Means + LR & 0.8782 & 0.7648 & 0.8421 & 0.6667 & 0.8387 & 0.9434 & 0.5333 \\
Fixed ridge recipe (157 gold) & 0.8932 & 0.7829 & 0.9412 & 0.6667 & 0.8000 & 0.9804 & 0.5263 \\
\bottomrule
\end{tabular}

\captionof{table}{Historical validation on 78 trials, fitted using 157 original gold training rows.
All rows use thresholding followed by argmax if empty. These are pre-refit diagnostic comparisons;
the final public ridge model subsequently uses all 235 visible labels. Positive support is
9 Safety, 4 Funding, 14 Efficacy, 50 Recruitment, and 10 Other.}
\label{tab:p1val}
\end{center}
\begin{center}\small
\begin{tabular}{lrrr}
\toprule
Evaluation & Fixed K-Means + LR & Local selection procedure & Fixed ridge recipe\\
\midrule
Repeated near-group OOF, 235 labels & 0.8877 & 0.8983 & 0.9026\\
Historical validation, 78 trials & 0.8782 & 0.8803 & 0.8932\\
Late train, 20 trials & 0.8417 & 0.8833 & 0.9000\\
Late combined, 56 trials & 0.8393 & 0.8869 & 0.8512\\
\bottomrule
\end{tabular}
\captionof{table}{Sample F1 under explicitly different procedures. The middle column evaluates
nested selection among local models, consensus students, and blends. The right column was
measured later for the already-selected public-winning ridge recipe; it is not an independent
estimate obtained before choosing that recipe. Temporal fits use 134 early train labels or
176 early train/validation labels respectively. The same labeled and unlabeled held groups
are excluded in OOF; temporal diagnostics retain earlier matching groups.}
\label{tab:p1eval}
\end{center}
}]
\subsection{What changed the conclusion}
Baseline F1 0.8654 omitted four empty predictions; the submission fallback raises it to 0.8782,
tying retained K-Means LR. Validation alone establishes no semi-supervised F1 gain.
Exact validation sets are 61/78 versus 64/78; test sets differ on 96/312 trials.
Earlier 157-label grouped OOF favored K-Means (0.8875 versus 0.8399), under different
grouping and populations from Table~\ref{tab:p1eval}.

Under the original empty-set policy, K-Means/GMM with LR $C=1$ score 0.8205/0.6667.
GMM produces no Safety, Funding, or Efficacy pseudo positives; K-Means produces no Funding
positives and only 39 Other positives among 1,014 masked rows. This label-quality limitation
motivates local models without inherited cluster targets.

\subsection{Grouped selection and calibration}
The 235 visible labels form 166 near-duplicate groups: connected components at raw Euclidean
distance $\leq10^{-4}$ across gold and masked rows. Held groups leave all fitting pools.
Four outer folds repeat with seeds 42, 137, and 2026; three inner group folds select models
and decisions, with further grouped calibration inside each fit. The 705 OOF predictions
repeat 235 trials, rather than represent independent observations.

When the local model and K-Means/LR control have complementary inner errors, local and
gold-only linear teachers can
supply label-wise soft pseudo targets. All six member scores must be at least 0.8 or at most
0.1; otherwise that label remains unknown. Gold weights remain one, pseudo weights are 0.1
with separate class-support caps, and only one student round is fitted. No outer fold selected
this student. Three of 12 folds selected blends; nine selected local models alone.

\subsection{Observed improvement and remaining errors}
Full-data selection chose ridge; 11 outer folds selected SVMs. Its public F1 is 0.88568,
up 0.03163 from K-Means LR. This aggregate feedback reveals no individual hidden labels.
The procedure's paired group-bootstrap 95\% gain interval over the inner-retuned control,
$[-0.00690,0.03494]$, includes zero.

The procedure improves Safety OOF F1 from 0.8228 to 0.9103, while Other declines from 0.5865
to 0.5761. Nine distinct Recruitment+Other trials are predicted as Recruitment alone in all
three repeats. The main improvement therefore does not resolve this rare co-occurrence.

\subsection{Feature ablation}
\label{sec:p1features}
With K-Means pseudo labels, gold/pseudo weights, LR $C=1$, and threshold 0.5 fixed,
we append 27 level-0 Reactome pathway indicators (training frequency $\geq5$) and two
availability flags to the 150 embeddings. Validation sample/macro F1 falls from
0.8205/0.6332 to 0.8120/0.5680. This controlled feature comparison uses the original
empty-set-allowing policy, unlike Table~\ref{tab:p1val}.
A separate TF-IDF view of criteria, diseases, and drugs supplies conservative Other additions;
it changes no decisions at its OOF-selected threshold: validation F1 remains 0.8782 and
late-56 F1 0.8393. This tests a targeted correction, not all text/embedding combinations.
